\documentclass[UTF8,a4paper,11pt]{ctexart}

\usepackage[a4paper,left=24mm,right=24mm,top=22mm,bottom=22mm]{geometry}
\usepackage{amsmath}
\usepackage{array}
\usepackage{booktabs}
\usepackage{enumitem}
\usepackage{fancyhdr}
\usepackage{longtable}
\usepackage{listings}
\usepackage{tabularx}
\usepackage{xcolor}
\usepackage[breaklinks]{hyperref}

\definecolor{linkblue}{RGB}{30,80,150}
\definecolor{codegray}{RGB}{245,246,248}
\definecolor{rulegray}{RGB}{90,90,90}
\definecolor{verified}{RGB}{0,105,70}
\definecolor{pending}{RGB}{170,80,0}
\definecolor{recommend}{RGB}{80,50,140}

\hypersetup{
  colorlinks=true,
  linkcolor=linkblue,
  urlcolor=linkblue,
  citecolor=linkblue,
  pdftitle={PlayCover PlayTools 直接渲染捕获开发方案},
  pdfauthor={代码调查报告}
}

\setlist{nosep,leftmargin=2em}
\setlength{\parindent}{2em}
\setlength{\parskip}{0.35em}
\setlength{\emergencystretch}{7em}
\sloppy
\setcounter{tocdepth}{2}

\lstset{
  basicstyle=\ttfamily\small,
  backgroundcolor=\color{codegray},
  frame=single,
  framerule=0.3pt,
  rulecolor=\color{rulegray},
  columns=fullflexible,
  breaklines=true,
  keepspaces=true,
  showstringspaces=false,
  xleftmargin=1em,
  xrightmargin=1em
}

\newcommand{\pathref}[2]{\path{#1}:#2}
\newcommand{\verified}{\textbf{\color{verified}[已证实]}}
\newcommand{\pending}{\textbf{\color{pending}[待验证]}}
\newcommand{\recommend}{\textbf{\color{recommend}[建议]}}
\newcommand{\must}{\textbf{必须}}
\newcommand{\should}{\textbf{应}}

\title{PlayCover / PlayTools 直接渲染捕获开发方案\\
  \large 面向《明日方舟》Unity Metal 渲染目标的 pre-compositor 图像接口}
\author{代码调查报告}
\date{2026 年 9 月 1 日}

\begin{document}
\pagestyle{plain}
\maketitle
\thispagestyle{empty}

\begin{abstract}
本方案针对 PlayCover 中现有窗口截图受逻辑尺寸、虚拟 native scale、Retina backing scale 和窗口装饰影响的问题，设计一条只适配 \texttt{com.hypergryph.arknights} 的直接渲染捕获链。目标是在被 PlayTools 注入的目标进程内，取得 Unity/Metal 提交给 CoreAnimation 之前的 color target，而不是再次从 macOS 窗口合成结果反向采样。

方案明确区分现状、设计和验证状态：当前固定 PlayTools revision 只有 ScreenCaptureKit/CGWindow 窗口采集；Arknights 的 UnityFramework 虽然包含 Unity Metal 和显示表面线索，但尚未证明具体运行时 drawable 属性或 hook ABI。推荐先部署不复制像素的运行时 probe，再在确认条件满足后，于原始 \texttt{presentDrawable:} 提交前编码 GPU copy，使用有界环形缓冲异步 readback，最后通过新增 MaaTools 命令传输带完整元数据的原始帧。旧 \texttt{SCRN}/\texttt{BGR\textbackslash{}1} 保持不变，失败时回退到原窗口采集。
\end{abstract}

\tableofcontents
\clearpage

\section{执行摘要}

\subsection{问题定义}

当前 MaaTools 的截图链路是：

\begin{lstlisting}
Unity / Metal render target
  -> PlayTools virtual screen and NSWindow
  -> ScreenCaptureKit or CGWindowListCreateImage
  -> title-bar crop and CGContext/vImage conversion
  -> MaaTools TCP response
\end{lstlisting}

该链路适合获得“用户在窗口中看到的图像”，但不适合作为稳定的渲染帧接口。PlayCover 的逻辑窗口尺寸、PlayTools 伪造的 UIKit native 尺寸、macOS 窗口的 backing pixels，以及图像采集 API 的实际返回像素并不一定相同。对 720p 做二次转换时，尺寸、裁剪、方向和通道语义容易发生漂移。

本方案把目标改为：

\begin{lstlisting}
Unity Metal render target
  -> in-process render tap
  -> GPU copy to bounded capture slots
  -> asynchronous readback
  -> MaaTools direct-render response
\end{lstlisting}

\subsection{结论与可行性等级}

\begin{center}
\small
\begin{tabularx}{\textwidth}{>{\raggedright\arraybackslash}p{0.30\textwidth} X}
\toprule
项目 & 判断 \\
\midrule
针对 Arknights 获取 Unity/Metal color target & \verified{} 技术上可行，但必须先完成运行时 probe。\\
绕过 NSWindow、ScreenCaptureKit、CGWindow 和 Retina backing scale & \verified{} 设计上可行。\\
仅修改 PlayCover 主程序 & 不可行；渲染对象位于注入后的目标应用进程。\\
复用当前 PlayTools 不做重建 & 不可行；固定 revision 没有 Metal capture 模块。\\
使用 \texttt{presentDrawable:} 直接 CPU 读像素 & 不可行；需要在 present 前 GPU copy，完成后异步读取。\\
准确输出真实 render target 尺寸 & 可行；尺寸必须来自 drawable/texture 元数据。\\
保证包含所有 CoreAnimation/UIKit overlay & 不保证；drawable 位于窗口合成之前。\\
保持旧 MAA 客户端可用 & 可行；保留旧命令，新增 direct-render 命令。\\
\bottomrule
\end{tabularx}
\end{center}

\subsection{推荐实施顺序}

\begin{enumerate}
  \item 固定目标应用 gate，并建立不修改像素的运行时诊断 probe。
  \item 观察 Unity 动态加载、\texttt{CAMetalLayer}、drawable 和实际 present selector。
  \item 记录 \texttt{framebufferOnly}、尺寸、格式、采样数和存储模式。
  \item 只复制一帧到自有 GPU/CPU 缓冲，验证同步和像素完整性。
  \item 新增 MaaTools direct-render 命令，保留旧截图命令。
  \item 生成 arm64/Catalyst 产物，经过 PlayCover 现有注入和重签名链部署。
  \item 执行窗口、旋转、resize、后台、性能、完整性和回退测试。
\end{enumerate}

\section{范围、依据与非目标}

\subsection{目标对象和版本边界}

本方案只针对以下本机测试对象：

\begin{center}
\small
\begin{tabularx}{\textwidth}{>{\raggedright\arraybackslash}p{0.32\textwidth} X}
\toprule
字段 & 已知值 \\
\midrule
Bundle identifier & \texttt{com.hypergryph.arknights} \\
应用版本 & 2.7.61，build 59 \\
主架构 & arm64 \\
引擎 & Unity 2021.3.39f1，IL2CPP \\
图形后端证据 & 直接链接 Metal/OpenGLES；运行时加载 AGXMetal 和 IOSurface \\
PlayCover & MaaTools 分支，HEAD \texttt{63cf605} \\
PlayTools 依赖 & \texttt{hguandl/PlayTools} MaaTools 分支，revision \texttt{7aa1977} \\
当前配置 & 内容目标 1280\(\times\)720、\texttt{customScaler=1}、MaaTools 端口 1717 \\
\bottomrule
\end{tabularx}
\end{center}

PlayTools 的固定 revision 和 Arknights 的 UnityFramework 都可能随版本变化。本文的 Unity 内部类名和 selector 只能作为当前版本的适配依据，不能作为通用 Unity API 或稳定 ABI。

\subsection{证据标签}

\begin{itemize}
  \item \verified{}：已经由仓库、固定依赖、IPA、已安装应用或已有运行时观察确认。
  \item \pending{}：需要新增 probe、构建产物或真实运行测试才能确认。
  \item \recommend{}：本方案建议采用的工程做法，不表示当前代码已经具备该能力。
\end{itemize}

\subsection{本方案包含的目标}

\begin{itemize}
  \item 在目标进程内获取 Unity/Metal pre-compositor color target。
  \item 让输出携带真实像素宽高、行距、格式、方向、帧号和状态。
  \item 消除窗口标题栏、AppKit frame 和 macOS backing scale 对原始帧尺寸的影响。
  \item 对 Arknights 做严格 gate，并在条件不满足时安全回退。
  \item 保持既有 MaaTools 命令的兼容性。
\end{itemize}

\subsection{明确的非目标}

\begin{itemize}
  \item 不把本方案扩展为所有 PlayCover 应用的通用渲染 hook。
  \item 不修改 UnityFramework 内部二进制或重新编译 Arknights 的游戏逻辑。
  \item 不获取 macOS 桌面、其他应用窗口或屏幕内容。
  \item 不绕过应用完整性、反作弊或服务端安全策略；完整性 SDK 只作为兼容性风险进行测试。
  \item 不在本阶段移除或改变现有窗口截图后端。
\end{itemize}

\section{现状基线与关键证据}

\subsection{PlayCover 主进程的职责}

PlayCover 主程序不持有游戏的 Metal command buffer。它负责将 PlayTools 安装到用户 framework 目录、向目标 Mach-O 添加 load command、复制 \texttt{AKInterface.bundle}、签名并启动应用。

\begin{itemize}
  \item \pathref{PlayCover/PlayCover/Model/PlayApp.swift}{64--116}：同步设置、检查 PlayTools、安装插件和启动目标应用。
  \item \pathref{PlayCover/PlayCover/Utils/PlayTools.swift}{46--91}：系统 framework 安装、Mach-O 注入和签名。
  \item \pathref{PlayCover/PlayCover/Utils/PlayTools.swift}{94--125}：将 \texttt{AKInterface.bundle} 复制进目标应用并重新签名。
  \item \pathref{PlayCover/PlayCover/Utils/PlayTools.swift}{150--187}：导出 IPA 时复制 \texttt{PlayTools.dylib} 并注入 \texttt{@executable\_path/Frameworks} 路径。
\end{itemize}

PlayCover 在 \pathref{PlayCover/PlayCover/Model/PlayApp.swift}{124--163} 清除 \texttt{DYLD\_*}、\texttt{METAL\_*} 和 \texttt{MTLCaptureEnabled} 等环境变量。这是启动环境管理，不是截图实现。因此 direct-render 不能依赖 Metal 调试环境变量注入，应该作为 PlayTools 内部代码随静态 load command 进入目标进程。

\subsection{现有 PlayTools 截图链}

固定 revision 的外部源码中，\texttt{MaaTools.swift} 的截图处理调用 \texttt{AKInterface.windowImage()}。\texttt{AKPlugin.swift} 再通过 macOS 14.4 以上的 ScreenCaptureKit 或旧版 CGWindow API 获取窗口图像，随后由 MaaTools 进行标题栏裁剪和栅格化。

\begin{center}
\small
\begin{tabularx}{\textwidth}{>{\raggedright\arraybackslash}p{0.28\textwidth} X}
\toprule
接口 & 当前语义 \\
\midrule
\texttt{SCRN} & 返回长度和 4 bytes/pixel 的原始 sRGB 缓冲；宽高由独立的 \texttt{SIZE} 提供。\\
\texttt{BGR\textbackslash{}1} & 返回裁剪后宽度、高度、长度和 3 bytes/pixel 缓冲；实现使用 RGBA8888 到 RGB888 转换，命令名不能单独证明通道顺序。\\
\texttt{SIZE} & 返回缓存的窗口内容尺寸，不是 Metal drawable 的可靠像素尺寸。\\
\bottomrule
\end{tabularx}
\end{center}

现有协议的握手和请求帧为：

\begin{lstlisting}
client  -> MAA\0
server  -> OKAY
request -> u16be payload_length || payload
payload[0..3] = command magic
\end{lstlisting}

\subsection{Arknights 的 Unity Metal 线索}

已安装 Arknights 的 UnityFramework 中可观察到以下静态线索：

\begin{lstlisting}
Unity 2021.3.39f1
GfxDeviceMetal
UnityDisplaySurfaceMTL
CAMetalLayer
nextDrawable
presentDrawable:
DisplayConnection.present
DisplayManager.present
UnityRenderingView.layerClass
UnityView.recreateRenderingSurface
MTLCommandBuffer
MTLTexture
setFramebufferOnly:
\end{lstlisting}

运行中的目标进程还加载了 \texttt{AGXMetalG17G.bundle}、IOSurface 和 IOAccelerator。上述事实支持 Metal 路径的研究，但不等于已经证明当前每一帧都由同一 selector 提交，也不等于已获得可复制的 drawable。具体调用频率、参数编码和对象生命周期必须由 probe 确认。

\subsection{尺寸边界}

当前 PlayCover 配置的内容目标为 1280\(\times\)720，窗口外部 bounds 观察到约为 1280\(\times\)748，显示器 backing scale 为 2.0。应用自己的 Unity 偏好中还出现宽高为 0、全屏和无边框的设置，意味着 Unity 可根据显示目标动态决定 render size。

方案中定义以下尺寸，避免混用：

\begin{center}
\small
\begin{tabularx}{\textwidth}{>{\ttfamily}p{0.18\textwidth} X}
\toprule
符号 & 含义 \\
\midrule
\texttt{W0,H0} & PlayCover/PlayTools 的逻辑窗口内容尺寸，接近 UIKit points。\\
\texttt{S} & PlayTools 伪造给应用的 native scale 或 custom scaler。\\
\texttt{B} & macOS 窗口实际 backing scale。\\
\texttt{Wr,Hr} & Unity/Metal 实际 render target 的 pixels，必须从 drawable/texture 读取。\\
\texttt{T} & 客户端可选的统一输出尺寸，例如 1280\(\times\)720。\\
\bottomrule
\end{tabularx}
\end{center}

direct-render 的源尺寸是 \texttt{Wr,Hr}，而不是 \texttt{W0,H0}、\texttt{W0\(\times\)B} 或窗口 frame。

\section{总体架构}

\subsection{进程和模块边界}

\begin{lstlisting}
PlayCover host process
  +-- AppSettings / install / injection / signing
  `-- launch target application
       |
       v
Arknights target process
  +-- PlayTools.dylib / framework
  |    +-- existing PlayScreen / PlayInput
  |    +-- existing MaaTools TCP server
  |    `-- new ArknightsMetalCapture
  +-- AKInterface.bundle
  `-- UnityFramework
       +-- UnityDisplaySurfaceMTL
       +-- CAMetalLayer
       `-- Metal command queue / command buffer
       |
       v
MaaCore client
  `-- direct-render command or legacy SCRN/BGR\1
\end{lstlisting}

\subsection{模块职责}

\begin{center}
\small
\begin{tabularx}{\textwidth}{>{\raggedright\arraybackslash}p{0.25\textwidth} X}
\toprule
模块 & 职责与边界 \\
\midrule
\texttt{ArknightsMetalCapture} & bundle gate、动态加载后安装 hook、probe、资源创建、GPU copy、slot 管理和异步 readback。不访问窗口，不执行网络发送。\\
\texttt{UnitySurfaceProbe} & 只记录 Unity surface/layer/drawable 的运行时元数据，首阶段不复制像素。\\
\texttt{MaaTools.swift} & 保存最新完成帧，处理 direct-render 命令和响应序列化；继续提供旧命令。\\
\texttt{AKPlugin.swift} & 保留现有窗口采集和完整窗口 fallback；不承担 Metal command buffer 同步。\\
PlayCover host & 默认负责现有设置、注入、复制和签名；只有增加用户可见开关时才扩展设置模型/UI。\\
MaaCore controller & 解析新响应、检查元数据和 stride，旧协议保持原解析逻辑。\\
\bottomrule
\end{tabularx}
\end{center}

\subsection{双层 hook 策略}

UnityFramework 不一定由 Arknights 主 Mach-O 直接链接，可能在启动后动态加载。因此不能只在 PlayLoader constructor 中立即查找 \texttt{DisplayConnection}。

\recommend{}采用两层策略：

\begin{enumerate}
  \item 在 PlayTools constructor 中只完成目标 gate、状态初始化和通用的低侵入观察准备。
  \item 监听 UnityFramework 动态加载，或使用受控的短周期延迟任务等待其进入地址空间。
  \item Unity 类出现后，以 Objective-C runtime 的 selector 和 method encoding 验证 \texttt{UnityRenderingView}、\texttt{DisplayConnection} 或 \texttt{DisplayManager} 适配点。
  \item 同时保留 \texttt{CAMetalLayer.nextDrawable} 诊断路径，作为版本专用 Unity hook 失效时的观察和兜底入口。
  \item 在所有 hook 中使用 \texttt{dispatch\_once}、递归保护和卸载/重建状态，避免重复安装。
\end{enumerate}

不建议使用 UnityFramework stripped 后的固定地址作为第一版方案。若未来必须使用地址特征，应把 UUID、版本和特征校验写成独立的适配层，并在不匹配时关闭 direct-render，而不是继续调用未知地址。

\section{阶段 M0：运行时 Probe}

\subsection{目的和原则}

M0 的目的不是拿到图像，而是回答 direct-render 是否具备实施条件。Probe 必须：

\begin{itemize}
  \item 只读对象属性，不修改 layer、drawable 或 command buffer；
  \item 不调用 \texttt{getBytes}、\texttt{waitUntilCompleted} 或同步网络；
  \item 对重复帧限频，只在对象、尺寸、格式或状态改变时记录；
  \item 只对目标 bundle 生效；
  \item 任何 selector、encoding 或对象类型不匹配都立即熔断。
\end{itemize}

\subsection{采集字段}

每次 layer/drawable 变化至少记录：

\begin{lstlisting}
process bundle identifier
UnityFramework UUID and version evidence
layer class and pointer identity
layer.contentsScale
layer.drawableSize.width / height
layer.pixelFormat
layer.framebufferOnly
 drawable class and texture class
texture.width / height
texture.pixelFormat
texture.sampleCount
texture.storageMode
texture.usage
command-buffer concrete class
present selector and Objective-C type encoding
thread identity and command-buffer status
\end{lstlisting}

日志不应记录账号、AppKey、网络请求正文或完整帧内容。指针地址仅用于当前进程内诊断，不应作为持久配置或协议字段。

\subsection{探测顺序}

\begin{enumerate}
  \item 确认 \texttt{Bundle.main.bundleIdentifier} 等于目标 bundle ID。
  \item 等待 UnityFramework 映像出现；记录 UUID 和路径，只作版本校验。
  \item 观察 \texttt{UnityRenderingView.layerClass} 或实际 view 的 layer 类型。
  \item 在 \texttt{nextDrawable} 返回后读取 layer、drawable 和 texture 的只读元数据。
  \item 观察 Unity surface 的 \texttt{present} 相关方法，确认是否每帧触发以及运行线程。
  \item 在启动、登录后场景、进入战斗、窗口 resize 和前后台切换时重复记录。
\end{enumerate}

\subsection{M0 判定表}

\begin{center}
\small
\begin{tabularx}{\textwidth}{>{\raggedright\arraybackslash}p{0.38\textwidth} X}
\toprule
观察结果 & 后续动作 \\
\midrule
稳定出现主 \texttt{CAMetalLayer}，drawable 尺寸非零，selector 可验证 & 进入 M1，先做单帧 GPU copy。\\
只有 Unity surface 方法，没有可观察 layer & 以 Unity surface 的 present 边界继续探测，不假设参数布局。\\
有 layer 但没有 present 调用 & 检查 present 变体或改为观察 command buffer；不在 nextDrawable 处读像素。\\
\texttt{framebufferOnly=NO} 且 texture 为单采样 & 可尝试 drawable blit。\\
\texttt{framebufferOnly=YES} & 进入 resolve/intermediate 分支，暂不承诺 drawable 可复制。\\
出现多个同尺寸 layer 或 layer 频繁重建 & 以生命周期和主 surface 识别为前置条件，未识别时关闭直采。\\
任何 selector encoding 不匹配或发生异常 & 记录原因并回退窗口采集。\\
\bottomrule
\end{tabularx}
\end{center}

\section{阶段 M1：Metal 捕获核心}

\subsection{捕获边界}

推荐的时间顺序是：

\begin{lstlisting}
Unity encodes render commands
  -> final resolve or blit
  -> capture copy encoded into the same command buffer
  -> original presentDrawable call
  -> command buffer commit and completion
  -> background queue reads capture slot
\end{lstlisting}

\texttt{presentDrawable:} 是一个提交边界候选，而不是 CPU 读回 API。hook 中应先确认 command buffer 仍处于 encoding 状态，然后编码对自有资源的 copy，最后调用原始 present 实现。不得在 hook 中等待 GPU，也不得在渲染线程同步发送 TCP。

\subsection{候选 hook 的优先级}

\begin{enumerate}
  \item Unity 版本专用的 \texttt{DisplayConnection.present} 或 \texttt{DisplayManager.present}：若 runtime encoding 可验证，优先使用，因为更接近 Unity 自己的 surface 管理。
  \item \texttt{CAMetalLayer.nextDrawable}：用于关联主 layer、记录 drawable 属性和捕获 resize；它本身不是最终帧完成点。
  \item \texttt{MTLCommandBuffer} concrete class 的 \texttt{presentDrawable:}：用于在 command buffer 中追加 copy；具体 class 由运行时 Metal 驱动决定。
  \item Unity 内部 stripped C++ 地址或指令特征：仅作为后续版本专用适配，不作为首选。
\end{enumerate}

\subsection{资源模型}

每个 capture slot 至少包含：

\begin{lstlisting}
slot state: free / encoding / inFlight / ready / dropped
width, height, bytesPerRow
pixel format and color-space tag
MTLBuffer or staging texture
frame ID and monotonic timestamp
orientation and rotation
MTLCommandBuffer completion status
\end{lstlisting}

\recommend{}使用两个或三个有界 slot。渲染线程只取得 free slot 并编码 copy；GPU 完成后由 completion handler 将 slot 标记为 ready；MaaTools 读取时只取最新 ready slot，并将过期 ready slot 丢弃。客户端慢时不应反向阻塞 render thread。

\subsection{GPU copy 和 readback}

优先使用以下路径：

\begin{enumerate}
  \item 创建与源尺寸一致的共享或 staging 资源；
  \item 按实际 row pitch 对齐 buffer；
  \item 在原 command buffer 中追加 texture-to-buffer 或 texture-to-texture copy；
  \item 调用原始 present；
  \item 在 completion handler 中确认 status/error；
  \item 将 ready slot 发布给专用串行 readback 队列；
  \item 读取前再次检查 slot 的 frame ID 和尺寸，防止 resize 期间误用旧资源。
\end{enumerate}

不允许以下实现：

\begin{itemize}
  \item 在 \texttt{presentDrawable:} 内调用同步 \texttt{waitUntilCompleted}；
  \item 直接对 framebuffer-only 或 private texture 调用未经验证的 CPU readback；
  \item 将 drawable 本身长期 retain 作为帧缓存；
  \item 每一帧额外调用 \texttt{nextDrawable}；
  \item 在 Metal 内部线程上进行阻塞式 socket 写入。
\end{itemize}

\subsection{framebuffer-only 分支}

\texttt{CAMetalLayer.framebufferOnly} 是第一项硬条件。其处理决策如下：

\begin{center}
\small
\begin{tabularx}{\textwidth}{>{\ttfamily}p{0.25\textwidth} X}
\toprule
状态 & 处理 \\
\midrule
\texttt{NO} & 检查 texture usage、storage、sampleCount 后尝试同 command buffer blit。\\
\texttt{YES}，可访问已 resolve intermediate target & 捕获 resolve 后的 intermediate target。\\
\texttt{YES}，无可访问中间目标 & 仅在 probe 环境测试创建前设为 NO；若改变渲染行为则关闭直采。\\
未知或对象已重建 & 不复制像素，等待下一次 probe；必要时回退窗口采集。\\
\bottomrule
\end{tabularx}
\end{center}

不能仅在 drawable 创建后修改 \texttt{framebufferOnly} 就假定已有 drawable 的使用限制已经改变。强制设为 NO 可能改变性能和 Unity 行为，必须作为 Arknights-only 的实验开关，不应默认对所有应用执行。

\subsection{MSAA、resolve 和 memoryless}

若 Unity 的渲染顺序为：

\begin{lstlisting}
MSAA render target -> resolve -> single-sample drawable -> present
\end{lstlisting}

在 drawable 可复制的前提下，捕获 drawable 通常能得到已 resolve 的最终颜色。以下情况必须由 probe 或 Metal validation 结果确认：

\begin{itemize}
  \item source texture 的 \texttt{sampleCount} 是否为 1；
  \item resolve 是否发生在当前 command buffer 的 copy 之前；
  \item render pass 是否使用 memoryless attachment；
  \item store action 是否保留 resolve 结果；
  \item 后处理、TAA、色调映射是否在另一个 command buffer 中追加；
  \item Unity 是否先绘制到 RenderTexture，再将最终结果 blit 到 drawable。
\end{itemize}

如果源是多采样纹理，必须先 resolve，不能把它按普通 2D 像素缓冲解释。如果 render pass 使用 discard 或 memoryless，过了该 render pass 后可能已没有可恢复的内容，此时 capture 点必须前移到 resolve 阶段。

\subsection{像素格式与色彩}

协议不得假定所有帧都是 4 bytes/pixel。第一版实现可优先支持实际 probe 得到的常见格式，例如 BGRA8Unorm、BGRA8Unorm\_sRGB、RGBA8Unorm 和 RGBA8Unorm\_sRGB；其他格式明确返回 unsupported 状态并回退。

每种支持格式应定义：

\begin{itemize}
  \item Metal pixel format 到协议枚举的映射；
  \item bytes per pixel 和 row pitch；
  \item R/G/B/A 通道排列；
  \item sRGB/linear/HDR 标识；
  \item 客户端是否需要转换为 BGR888 或 RGBA8888。
\end{itemize}

Arknights 的 boot 配置包含 \texttt{hdr-display-enabled=0}，这使 SDR 8-bit 路径更可能成立，但不能代替运行时检查。若出现 RGBA16Float 等 HDR 格式，应先转成明确的 SDR 输出，或返回不支持状态；不能把 float16 数据伪装成 8-bit 数据。

颜色通道必须使用红、绿、蓝色块探针验证。命令名 \texttt{BGR\textbackslash{}1} 不足以证明实际字节顺序，因为现有实现的 vImage 调用语义与命令名并不完全一致。

\subsection{方向和旋转}

Metal texture 通常不携带用户界面方向语义，输出方向需要单独记录。每一帧元数据应包含：

\begin{lstlisting}
orientation: portrait / landscape-left / landscape-right / unknown
rotation: 0 / 90 / 180 / 270
mirrored: false or true
content origin: top-left or bottom-left
\end{lstlisting}

MaaTools 和 MaaCore 应把“像素内存方向”和“触控坐标方向”分开处理。不要把 PlayCover 的 \texttt{displayRotation} 直接当作 Metal texture 的旋转；必须用已知角落文字或测试图验证。

\subsection{overlay 语义}

direct-render 返回的是 Unity/Metal drawable，不是 macOS 最终合成画面。它通常不包含：

\begin{itemize}
  \item macOS 标题栏、边框和阴影；
  \item 位于 Unity view 上方的独立 UIKit/CoreAnimation layer；
  \item 系统键盘或独立窗口；
  \item PlayTools 创建的非 drawable 浮层。
\end{itemize}

已经绘制进 Unity 最终 render pass 的游戏 UI 会包含在内。方案应将该语义写入协议能力说明，并保留现有窗口采集作为“完整窗口”后端。

\section{阶段 M2：MaaTools 协议和客户端}

\subsection{兼容原则}

旧客户端不认识未知 magic，且旧命令已经被使用。因此：

\begin{itemize}
  \item 不改变 \texttt{SCRN}、\texttt{BGR\textbackslash{}1}、\texttt{SIZE} 的已有布局和语义；
  \item 新增独立 direct-render 命令，例如 \texttt{RND0}；
  \item 只有支持该命令的客户端才启用直采；
  \item direct-render 失败时，客户端可以继续请求旧命令；
  \item \texttt{VERN=3} 不应被偷偷改成另一种旧命令语义；如需能力发现，增加显式能力查询或在新响应中返回能力状态。
\end{itemize}

\subsection{建议的 RND0 响应}

请求可以沿用现有握手和 16-bit 大端长度帧：

\begin{lstlisting}
connect:
    client  -> MAA\0
    server  -> OKAY

request:
    u16be payload_length || "RND0" || request_flags
\end{lstlisting}

建议响应使用固定大端序头部，再跟随原始字节。具体版本应在实现前冻结并写入单元测试：

\begin{center}
\small
\begin{tabularx}{\textwidth}{>{\ttfamily}p{0.25\textwidth} X}
\toprule
字段 & 含义 \\
\midrule
\texttt{magic} & 响应 magic，例如 \texttt{RND0}。\\
\texttt{version:u16} & direct-render 响应版本，与旧 MaaTools 协议版本分开。\\
\texttt{status:u16} & ready、no-frame、unsupported-format、not-target、temporarily-disabled 等。\\
\texttt{frameID:u64} & 单调递增的帧序号。\\
\texttt{timestamp:u64} & 单调时钟纳秒或约定的时间单位。\\
\texttt{width:u32} & render target 像素宽度。\\
\texttt{height:u32} & render target 像素高度。\\
\texttt{bytesPerRow:u32} & 实际行距，不能默认等于 width\(\times\)BPP。\\
\texttt{pixelFormat:u32} & 协议枚举，不直接暴露平台私有数值。\\
\texttt{orientation:u16} & 方向枚举。\\
\texttt{rotation:u16} & 旋转枚举。\\
\texttt{dataLength:u32} & 后续像素字节数。\\
\texttt{data} & 按上述格式和 row pitch 编码的像素数据。\\
\bottomrule
\end{tabularx}
\end{center}

若状态不是 ready，响应仍应带有可解析的 status 和错误原因，不能让客户端把空数据当作有效的零尺寸帧。长度计算必须使用 64-bit 中间值并在写入 u32 前检查上限。

\subsection{服务器状态机}

\begin{lstlisting}
Disabled
  -> TargetVerified
  -> ProbeReady
  -> CaptureArmed
  -> FrameInFlight
  -> FrameReady
  -> ClientRead

任何 selector mismatch / GPU error / resize race
  -> TemporarilyDisabled
  -> legacy window capture fallback
\end{lstlisting}

MaaTools 不应为了响应一次请求临时创建 drawable，也不应等待下一帧无限阻塞连接。没有 ready 帧时返回明确状态；客户端可短暂重试或请求旧后端。

\subsection{MaaCore 适配}

若采用新命令，MaaCore 的 PlayTools controller 需要新增响应解析和能力选择。已知的外部客户端控制器位置为：

\begin{lstlisting}
maa/src/MaaCore/Controller/PlayToolsController.h
maa/src/MaaCore/Controller/PlayToolsController.cpp
\end{lstlisting}

客户端必须：

\begin{enumerate}
  \item 按大端序读取固定头部；
  \item 校验 status、width、height、bytesPerRow 和 dataLength；
  \item 检查 \texttt{dataLength} 是否覆盖所有行距；
  \item 根据 pixelFormat 解码，而不是由命令名猜通道；
  \item 根据 orientation/rotation 做显示或识别坐标变换；
  \item 在 direct-render 不支持时自动使用旧 \texttt{SCRN} 或 \texttt{BGR\textbackslash{}1}。
\end{enumerate}

\section{阶段 M3：配置、gate 与回退}

\subsection{Arknights gate}

第一版必须在目标进程内执行硬 gate：

\begin{lstlisting}
Bundle.main.bundleIdentifier == "com.hypergryph.arknights"
\end{lstlisting}

可选的第二层校验包括 UnityFramework UUID、Unity 版本字符串和目标 framework 路径。任何校验失败都关闭 direct-render，不影响 PlayTools 的普通输入和旧截图服务。

\subsection{开关策略}

\recommend{}第一版不要立即增加复杂的 PlayCover UI。可以先将 direct-render 作为 MaaTools 内部的 Arknights-only 能力：

\begin{itemize}
  \item 旧 \texttt{maaTools} 仍决定 MaaTools listener 是否启动；
  \item 新能力只有在 Arknights gate 和 probe 同时成功后自动 armed；
  \item 客户端通过新命令请求，未请求时不做像素 copy；
  \item 发生失败时只记录一次原因并回退。
\end{itemize}

若后续需要用户可见的“Metal direct capture”开关，再同步修改：

\begin{itemize}
  \item \pathref{PlayCover/PlayCover/Model/AppSettings.swift}{10--56}：新增可编码设置及默认值；
  \item \pathref{PlayCover/PlayCover/Views/AppSettingsView.swift}{576--586}：新增开关和说明；
  \item 外部 PlayTools 的 \texttt{PlaySettings.swift}：读取同一个 plist 字段；
  \item 本地化资源：解释该模式不包含独立 CoreAnimation overlay。
\end{itemize}

\subsection{回退条件}

以下任一条件出现时，direct-render 应暂时禁用并回到 \texttt{AKInterface.windowImage()}：

\begin{itemize}
  \item bundle、Unity UUID 或 selector encoding 不匹配；
  \item drawable 尺寸为零或尺寸变化尚未完成资源重建；
  \item \texttt{framebufferOnly} 阻止复制且没有可用 resolve target；
  \item sampleCount、pixelFormat 或 storage mode 不在支持矩阵；
  \item command buffer status/error 表示 copy 或 present 失败；
  \item ring buffer 连续耗尽、GPU 延迟超过阈值或读取数据校验失败；
  \item 应用进入后台、Unity layer 被重建或目标进程退出。
\end{itemize}

回退不应删除 hook 或修改 Unity 层；只需停止 armed 状态，释放 capture slot，并让旧窗口采集继续工作。

\subsection{网络和权限边界}

direct-render 访问的是目标进程自己的 Metal 资源，不调用 ScreenCaptureKit，因此理论上不需要 Screen Recording TCC。现有 MaaTools 的网络服务仍需单独处理：

\begin{itemize}
  \item listener 明确绑定 127.0.0.1，而不是依赖标题中的 \texttt{localhost}；
  \item 对新命令设置请求长度、连接超时和响应大小限制；
  \item 截图能力与 TOUCH/TERM 控制能力分离授权；
  \item 不将像素帧写入持久日志；
  \item 端口冲突或非预期来源时返回可诊断错误。
\end{itemize}

\section{阶段 M4：构建、注入和部署}

\subsection{需要修改的外部工程}

当前 PlayTools 源码不在 PlayCover checkout 中，固定 revision 的实际修改应发生在外部 PlayTools 工程：

\begin{center}
\small
\begin{tabularx}{\textwidth}{>{\raggedright\arraybackslash}p{0.35\textwidth} X}
\toprule
文件/target & 计划变更 \\
\midrule
PlayTools target & 新增 Metal capture 源文件，加入 Metal/QuartzCore 依赖，保留现有 MaaTools。\\
\texttt{MaaTools.swift} & 新命令、最新帧状态、序列化、fallback 和 listener 生命周期。\\
\texttt{AKPlugin.swift} & 继续提供窗口采集；若 probe 需要 layer 发现，可增加最小适配桥接，但不复制渲染帧。\\
\texttt{PlayCover.swift}/\texttt{PlayLoader.m} & 目标 gate、动态加载后初始化或 probe 启动入口。\\
\path{PlayTools.xcodeproj/project.pbxproj} & 新源文件、framework link、编译条件和 arm64/Catalyst target。\\
\texttt{PlaySettings.swift} & 只有新增 direct-render 配置字段时才修改。\\
\bottomrule
\end{tabularx}
\end{center}

\subsection{PlayCover host 的变更边界}

如果仍复用用户级 \texttt{PlayTools.framework} 路径，\pathref{PlayCover/PlayCover/Utils/PlayTools.swift}{46--187} 的注入、复制和签名逻辑原则上可以复用。当前 host 不需要知道 Metal texture，也不需要访问目标进程内存。

只有以下情况需要修改 PlayCover host：

\begin{itemize}
  \item 新增独立 dylib，导致需要额外复制或注入路径；
  \item 新增用户设置开关，需要修改设置模型和 UI；
  \item 新 framework 资源布局改变，导致 \texttt{installPluginInIPA} 复制规则改变；
  \item 新协议能力需要 host 在启动前传递非 plist 参数。
\end{itemize}

不建议通过 \texttt{DYLD\_INSERT\_LIBRARIES} 作为正式方案，因为 PlayCover 会清理 DYLD 环境变量，而且环境变量路径会增加首帧竞态和签名复杂度。

\subsection{Carthage 和 Catalyst 产物}

PlayCover 的 \pathref{PlayCover/PlayCover.xcodeproj/project.pbxproj}{578--597} 构建脚本当前：

\begin{enumerate}
  \item 从 \texttt{Carthage/Build/PlayTools.xcframework/ios-arm64} 读取 framework；
  \item 复制 \texttt{AKInterface.bundle}；
  \item 用 \texttt{vtool} 设置 maccatalyst build version；
  \item 重新组织 deep framework 目录；
  \item 对插件和 framework 签名。
\end{enumerate}

开发顺序应为：

\begin{lstlisting}
external PlayTools source
  -> build arm64 iOS/Catalyst-compatible framework
  -> verify Metal.framework and symbols/selectors
  -> Carthage/Build/PlayTools.xcframework
  -> PlayCover copy/vtool/rebundle
  -> codesign nested plugin and framework
  -> inject target Mach-O
  -> codesign target application
  -> launch Arknights and inspect logs
\end{lstlisting}

只有在确实改变了 PlayTools 来源 revision 后，才更新 [L] \texttt{Cartfile.resolved}；不能仅为生成本地测试产物而把未审计的依赖 revision 写入发布配置。

\subsection{签名和部署检查}

每次部署必须检查：

\begin{itemize}
  \item PlayTools、AKInterface 和目标应用均为 arm64；
  \item Metal/QuartzCore load command 在最终产物中存在；
  \item bundle gate 代码进入目标进程，而不是只进入 PlayCover host；
  \item nested bundle 先签、外层 framework 后签、目标 app 最后签；
  \item 复制资源前清理会导致 codesign 失败的扩展属性；
  \item 目标 Mach-O 的 PlayTools load command 仍然存在且路径有效；
  \item 替换现有用户级 framework 后重新启动应用，不复用旧进程。
\end{itemize}

\section{M5：测试与验收矩阵}

\subsection{M0 Probe 测试}

\begin{center}
\footnotesize
\begin{longtable}{>{\raggedright\arraybackslash}p{0.24\textwidth} >{\raggedright\arraybackslash}p{0.34\textwidth} >{\raggedright\arraybackslash}p{0.28\textwidth}}
\toprule
场景 & 验证内容 & 通过标准 \\
\midrule
\endfirsthead
\toprule
场景 & 验证内容 & 通过标准 \\
\midrule
\endhead
启动首帧 & UnityFramework 动态加载和 hook 安装时序 & 不崩溃，能在首个稳定 frame 前观察对象 \\
普通主界面 & layer、drawable、present 关联 & 主 layer 的属性稳定，present 连续触发 \\
进入战斗 & 渲染 target 和后处理变化 & 尺寸/格式变化被记录，不使用旧 slot \\
窗口 resize & drawableSize、layer 重建 & 资源重建，无旧尺寸帧混入 \\
前后台切换 & command buffer 状态和 layer 生命周期 & 暂停/恢复可识别，失败可回退 \\
非目标应用 & bundle gate & 不安装或不执行 Arknights capture \\
selector 不匹配 & encoding 和版本校验 & 禁用直采并记录一次原因 \\
\bottomrule
\end{longtable}
\end{center}

\subsection{M1 GPU copy 测试}

\begin{itemize}
  \item 单帧 copy 后检查非零长度、尺寸和 frame ID；
  \item 检查 GPU validation 是否报告 framebuffer-only、usage 或 encoder 错误；
  \item 检查 completion handler 后再读取，禁止提前读取；
  \item 连续运行至少数百帧，确认没有 drawable starvation、死锁或明显卡顿；
  \item 使用色块/角落文字验证通道顺序、行距、上下方向和旋转；
  \item 在 resize 后确认每个 slot 的容量和 bytesPerRow 已更新。
\end{itemize}

\subsection{协议测试}

\begin{center}
\small
\begin{tabularx}{\textwidth}{>{\raggedright\arraybackslash}p{0.31\textwidth} X}
\toprule
测试组 & 内容 \\
\midrule
握手与端序 & \texttt{MAA\textbackslash{}0/OKAY}、u16be 请求长度、所有 u32/u64 大端字段。\\
边界长度 & 零帧、最大允许尺寸、dataLength 溢出、半包、断开和重复请求。\\
格式校验 & width、height、stride、pixelFormat、dataLength 一致性。\\
能力和失败 & not-target、no-frame、unsupported-format、temporarily-disabled。\\
旧兼容 & 旧客户端继续使用 SCRN/BGR\textbackslash{}1；新客户端可回退旧命令。\\
并发 & 多次请求、慢客户端、断线重连、帧丢弃和只保留最新帧。\\
控制隔离 & direct-render 不扩大 TOUCH/TERM 权限，listener 只绑定 loopback。\\
\bottomrule
\end{tabularx}
\end{center}

\subsection{显示和几何测试}

\begin{itemize}
  \item 逻辑窗口 1280\(\times\)720，macOS backing scale 1 和 2；
  \item 窗口化、全屏、borderless 和标题栏显示；
  \item 16:9、4:3 和可调整窗口；
  \item displayRotation 0、90、180、270 和 inverse screen values；
  \item resize 过程中连续请求帧；
  \item Unity UI、可能的 UIKit overlay 与窗口采集结果对照；
  \item 统一输出 1280\(\times\)720 时分别验证 aspect-fit、aspect-fill 和黑边 validRect。
\end{itemize}

\subsection{性能和稳定性验收}

建议记录以下指标，而不是只观察“能否截图”：

\begin{lstlisting}
capture request latency
GPU copy duration if available
command-buffer completion latency
readback queue depth
number of dropped frames
number of fallback transitions
render thread frame-time delta
resident memory of capture slots
TCP response throughput
\end{lstlisting}

第一版验收条件：

\begin{enumerate}
  \item 目标应用在 probe 和 direct-render 开关打开时可正常启动和退出。
  \item 不因 capture 产生 Metal validation error、死锁、drawable starvation 或可重复崩溃。
  \item 输出尺寸来自真实 render target，且 stride、格式和数据长度可自洽验证。
  \item 客户端能正确处理方向和通道，旧客户端不受影响。
  \item resize、后台、Unity layer 重建和不支持格式都能回退。
  \item 未请求 direct-render 时不进行不必要的每帧 GPU copy。
  \item 非 Arknights 应用不安装、不执行该专用 capture 路径。
\end{enumerate}

\section{风险、取舍与回滚}

\subsection{技术风险}

\begin{center}
\small
\begin{tabularx}{\textwidth}{>{\raggedright\arraybackslash}p{0.28\textwidth} >{\raggedright\arraybackslash}p{0.32\textwidth} X}
\toprule
风险 & 影响 & 缓解 \\
\midrule
UnityFramework 动态加载 & constructor 查找不到 Unity 类或错过首帧 & 动态加载后初始化，加通用 layer probe，版本不匹配则关闭。\\
私有 concrete Metal class & selector/实现随系统 GPU 变化 & 运行时 class/encoding 校验，避免固定地址，保留 fallback。\\
framebuffer-only & drawable 不可作为 copy source & 先测实际值，优先捕获 resolve target，强制 NO 只做实验分支。\\
MSAA/memoryless & 过晚捕获得不到有效颜色 & 确认 resolve/store 时序，必要时前移到 intermediate target。\\
GPU 同步 & 卡顿、旧帧、资源复用冲突 & 同 command buffer copy、有界 slot、completion 后 readback。\\
overlay 差异 & direct frame 与用户所见画面不同 & 在协议中声明语义，保留窗口采集后端。\\
Unity/系统更新 & 私有 ABI 失效 & UUID/版本 gate、probe、失败熔断和回退。\\
\bottomrule
\end{tabularx}
\end{center}

\subsection{安全和合规风险}

目标应用包含第三方完整性/安全组件。即使本方案只使用公开 Metal API，进程内 hook 仍可能改变调用路径和内存行为。测试应使用用户明确授权的本机应用、副本或测试账号；不应把本方案用于绕过完整性检查、修改游戏逻辑或扩大控制权限。

现有 MaaTools listener 还存在网络暴露和 TOUCH/TERM 混用风险。direct-render 不是解决这些问题的理由，部署时应将 loopback、认证、超时和控制命令隔离作为独立验收项。

\subsection{回滚方案}

回滚必须是可逆的：

\begin{enumerate}
  \item 将 direct-render armed 状态置为 false；
  \item 停止新命令的像素 copy 和 readback；
  \item 释放 ring buffer，但不销毁现有 PlayScreen/PlayInput；
  \item 让 \texttt{AKInterface.windowImage()} 继续服务 \texttt{SCRN}/\texttt{BGR\textbackslash{}1}；
  \item 若新 framework 影响启动，则恢复上一份已验证的 PlayTools framework 和目标应用签名产物。
\end{enumerate}

不要使用不可恢复的二进制覆盖、强制 reset 或删除用户应用数据来处理失败部署。每次实验前应保留旧 framework hash、旧目标应用包和配置备份。

\section{里程碑与交付物}

\begin{center}
\small
\begin{longtable}{>{\raggedright\arraybackslash}p{0.13\textwidth} >{\raggedright\arraybackslash}p{0.39\textwidth} >{\raggedright\arraybackslash}p{0.34\textwidth}}
\toprule
里程碑 & 工作内容 & 退出条件 \\
\midrule
\endfirsthead
\toprule
里程碑 & 工作内容 & 退出条件 \\
\midrule
\endhead
M0 Probe & Arknights gate、动态加载观察、layer/drawable/present 元数据日志 & 真实主 layer 和 present 路径可解释，或明确记录阻塞条件 \\
M1 单帧 copy & 资源创建、framebuffer-only 分支、同 command buffer blit、completion readback & 单帧数据可校验，无 GPU 错误和渲染线程阻塞 \\
M2 协议接入 & RND0 头部冻结、服务端状态机、MaaCore 解析和 fallback & 新旧客户端互不破坏，错误状态可解析 \\
M3 构建部署 & PlayTools framework、AKInterface、Carthage、Catalyst、注入和嵌套签名 & 全新启动应用并加载正确 UUID/版本的 PlayTools \\
M4 回归 & resize、方向、格式、性能、后台、端口和完整性风险 & 验收矩阵通过，失败可自动回退 \\
M5 可选产品化 & host 设置/UI、能力说明、日志和发布脚本 & 用户可控开关、升级策略和回滚包齐全 \\
\bottomrule
\end{longtable}
\end{center}

\section{实施文件清单}

\subsection{第一阶段实际需要新增或修改的文件}

\begin{center}
\small
\begin{tabularx}{\textwidth}{>{\raggedright\arraybackslash}p{0.44\textwidth} X}
\toprule
文件 & 计划 \\
\midrule
外部 PlayTools 的新 Metal capture 源文件 & 新增；负责 probe、hook、GPU copy、slot 和 readback。\\
外部 PlayTools \texttt{MaaTools.swift} & 修改；增加 direct-render command、元数据和 fallback。\\
外部 PlayTools \texttt{PlayCover.swift}/\texttt{PlayLoader.m} & 修改或增加入口；处理动态加载后初始化。\\
外部 PlayTools \path{PlayTools.xcodeproj/project.pbxproj} & 修改；加入源码和 Metal framework link。\\
外部 MaaCore \texttt{PlayToolsController.h/.cpp} & 若采用 RND0，则修改客户端解析和能力回退。\\
\pathref{PlayCover/PlayCover/Model/AppSettings.swift}{10--56} & 第一版可不改；只有需要 host 开关时修改。\\
\pathref{PlayCover/PlayCover/Views/AppSettingsView.swift}{576--586} & 第一版可不改；产品化开关时修改。\\
\pathref{PlayCover/PlayCover/Utils/PlayTools.swift}{46--187} & 同一 framework 路径复用时不改；独立 dylib/资源布局变化时才改。\\
\pathref{PlayCover/PlayCover/Model/PlayApp.swift}{124--163} & 不用环境变量方案时不改；禁止依赖其已清理的 capture 环境变量。\\
\bottomrule
\end{tabularx}
\end{center}

\subsection{本次文档交付不修改的内容}

本次任务只生成本方案的 LaTeX 和 PDF，不执行上述 PlayTools 实现。不会修改：

\begin{itemize}
  \item PlayCover 应用源码；
  \item 外部 PlayTools 源码或二进制；
  \item MaaCore 源码；
  \item 用户提供的 IPA 和已安装 Arknights；
  \item PlayCover 的依赖锁定和签名配置。
\end{itemize}

\section{诊断日志和协议验收模板}

\subsection{建议的 probe 日志}

\begin{lstlisting}
[MetalCapture] target=com.hypergryph.arknights
[MetalCapture] UnityFramework uuid=<validated>
[MetalCapture] layer=<class> drawable=1280x720 scale=1.0
[MetalCapture] pixelFormat=<enum> framebufferOnly=NO sampleCount=1
[MetalCapture] texture storage=<enum> usage=<mask> rowPitch=<value>
[MetalCapture] present selector=<validated> thread=<id>
[MetalCapture] capture armed; slots=3
[MetalCapture] frame=42 bytes=... status=ready
\end{lstlisting}

实际日志中不应把尖括号内容当作固定值，也不应写入账号或应用网络数据。对同一个进程，layer 地址和 frame ID 只用于诊断；重启后不得依赖地址稳定性。

\subsection{协议不变量}

实现和客户端测试应固定以下不变量：

\begin{enumerate}
  \item 所有整数端序在协议文档中明确为大端序。
  \item \texttt{dataLength} 不超过服务端允许的最大响应大小。
  \item \texttt{bytesPerRow} 至少覆盖一行所需像素字节，并能覆盖全部数据长度。
  \item width、height、pixelFormat、stride 和 dataLength 相互一致。
  \item frame ID 单调递增；丢帧不要求连续，但不得回退。
  \item ready 帧在完成回调之前不可被客户端读取。
  \item status 非 ready 时，客户端不能把 data 当作有效图像。
  \item direct-render 失败不改变旧命令的返回布局。
\end{enumerate}

\section{最终建议}

本项目最稳妥的实施路线不是直接把现有 \texttt{SCRN} 改成 Metal 读取，而是新增一个具备明确语义的 direct-render 后端：

\begin{lstlisting}
probe -> validate -> arm -> GPU copy -> async readback -> RND0
             |                                      |
             +---------- failure / unsupported -----+
                              |
                     legacy window capture
\end{lstlisting}

在实际编写 hook 之前，M0 probe 是不可省略的验收门槛。特别是 \texttt{framebufferOnly}、resolve 时序、动态加载和 concrete command-buffer class，无法由静态字符串或 \texttt{nm} 输出替代。

只要当前实例满足可复制 drawable 或可访问的 resolve target，Arknights-only 方案就具备继续开发的条件；如果两个条件都不满足，则应保留现有窗口采集，而不是通过同步读回或未经验证的私有偏移强行实现。这样可以同时解决尺寸一致性问题、控制性能风险，并使失败路径保持可逆。

\appendix
\section{证据索引}

\small
\begin{longtable}{>{\raggedright\arraybackslash}p{0.11\textwidth} >{\raggedright\arraybackslash}p{0.38\textwidth} >{\raggedright\arraybackslash}p{0.39\textwidth}}
\toprule
标签 & 文件或对象 & 支持的事实 \\
\midrule
\endfirsthead
\toprule
标签 & 文件或对象 & 支持的事实 \\
\midrule
\endhead
\textbf{[L]} & \pathref{PlayCover/PlayCover/Model/PlayApp.swift}{64--116} & PlayCover 同步设置、检查 PlayTools、安装插件并启动目标应用。\\
\textbf{[L]} & \pathref{PlayCover/PlayCover/Model/PlayApp.swift}{124--163} & 清除 DYLD 和 Metal debug/capture 环境变量。\\
\textbf{[L]} & \pathref{PlayCover/PlayCover/Utils/PlayTools.swift}{46--187} & framework 复制、Mach-O 注入、插件复制和签名。\\
\textbf{[L]} & \pathref{PlayCover/PlayCover/Model/AppSettings.swift}{10--32} & MaaTools、端口、窗口尺寸和 custom scaler 设置字段。\\
\textbf{[L]} & \pathref{PlayCover/PlayCover/Views/AppSettingsView.swift}{576--586} & MaaTools 开关和端口 UI。\\
\textbf{[L]} & \pathref{PlayCover/PlayCover.xcodeproj/project.pbxproj}{578--635} & Carthage framework 复制、Catalyst 转换、重打包和签名脚本。\\
\textbf{[L]} & \pathref{PlayCover/Cartfile}{1}、\pathref{PlayCover/Cartfile.resolved}{1} & PlayTools 来源和固定 revision。\\
\textbf{[X]} & PlayTools@7aa1977：\texttt{MaaTools.swift}、\texttt{AKPlugin.swift} & 现有截图经 AKInterface.windowImage，再走 ScreenCaptureKit/CGWindow；没有 direct Metal capture。\\
\textbf{[X]} & PlayTools@7aa1977：\texttt{PlayCover.swift}、\texttt{PlayLoader.m} & 目标进程初始化顺序和 constructor 入口。\\
\textbf{[R]} & 已安装 Arknights UnityFramework & Unity Metal、CAMetalLayer、Unity surface 和 present 相关静态线索。\\
\textbf{[R]} & 已运行目标进程 & PlayTools、AKInterface、UnityFramework、AGXMetal、IOSurface 和 IOAccelerator 已加载。\\
\textbf{[I]} & Unity 私有 surface/command-buffer 方法 & 可作为当前版本 hook 候选，但必须通过 selector encoding 和运行时调用验证。\\
\bottomrule
\end{longtable}

\section{当前状态声明}

本文是开发方案，不是功能实现报告。文中所有 M0 之后的 hook、GPU copy、RND0 协议、MaaCore 修改、构建产物和性能数字均为待实施或待验证项目；当前已完成的是 PlayCover、固定 PlayTools revision、Arknights IPA/安装包和运行时模块的调查。

\end{document}
